\subsection{Blowup of Smooth Solutions}\label{sec:conservation-smooth-blowup}

As we noted above, the main purpose of this chapter is to study PDEs with discontinuous solutions. We are now prepared to show how discontinuous solutions of conservation laws can develop from continuous ones.

\subsubsection{Single Conservation Laws}\label{subsec:scalar-conservation-blowup}

We consider a single conservation law of the form
\begin{equation}\label{eq:scalar-conservation-law-quasilinear}
u_t+f'(u)u_x=0.
\end{equation}
Here $f$ is assumed sufficiently smooth. Characteristic curves for \eqref{eq:scalar-conservation-law-quasilinear} must satisfy
\begin{equation}\label{eq:scalar-conservation-characteristic-ode}
\hat x'(t)=f'(u(\hat x(t),t)).
\end{equation}


As a result of this relation we get the following very strong result for single conservation laws.

\begin{Theorem}
\label{thm:scalar-constant-characteristics}
Any $C^1$ solution of the single conservation law \eqref{eq:scalar-conservation-law-quasilinear} is constant along characteristics. Accordingly, characteristic curves for \eqref{eq:scalar-conservation-law-quasilinear} are straight lines.
\end{Theorem}

\begin{Proof}
Using \eqref{eq:scalar-conservation-law-quasilinear} and \eqref{eq:scalar-conservation-characteristic-ode}, we get
\begin{equation}\label{eq:scalar-solution-constant-along-characteristics}
\dfrac{\mathrm{d}}{\mathrm{d}t}u(\hat x(t),t)=u_x\hat x'+u_t=u_xf'(u)+u_t=0.
\end{equation}
Thus
\begin{equation}\label{eq:scalar-characteristic-constant-value}
u(\hat x(t),t)=C
\end{equation}
where $C$ is a constant. So \eqref{eq:scalar-conservation-characteristic-ode} implies
\begin{equation}\label{eq:scalar-straight-characteristic-line}
\hat x(t)=kt+\hat x(0),
\end{equation}
where $k$ is the constant $k:=f'(C)$.
\end{Proof}

\begin{figure}[H]
\centering
\includegraphics[width=.78\linewidth]{figures/characteristic-construction-source.png}
\caption{Defining a solution by characteristics.}\label{fig:characteristic-construction}
\end{figure}

To see what this implies in general about solutions of the Cauchy problem let us focus on Burgers' equation
\begin{equation}\label{eq:burgers-characteristic-cauchy-equation}
u_t+uu_x=0
\end{equation}
with the initial condition
\begin{equation}\label{eq:burgers-characteristic-cauchy-data}
u(x,0)=u_0(x).
\end{equation}
Note that the equation for characteristics reduces to
\begin{equation}\label{eq:burgers-characteristic-slope}
\hat x'(t)=u(\hat x(t),t).
\end{equation}


Thus, the initial data give us the slopes of the characteristic rays emanating from the $x$ axis. For certain initial data this gives us a method for "solving" the Cauchy problem. We simply go along the $x$ axis, drawing characteristic rays with slope depending on the initial data, and let the solution take the value of the corresponding initial data along the characteristic (cf. Figure~\ref{fig:characteristic-construction}).

Unfortunately, some simple examples of discontinuous initial data show us just how easily the procedure falls apart. In Figure~\ref{fig:characteristic-gap} we see that for an initial condition corresponding to a step function, there is a region that is untouched by any characteristics from the initial data; the procedure above does not identify a solution in this region. As we shall see below, in this case we will be able to identify a continuous solution called a rarefaction or fan wave.

However, in Figure~\ref{fig:characteristic-overlap} we have a more difficult problem. For a decreasing step function the characteristics overlap. Since our solution cannot be multivalued, we must conclude (in light of Theorem~\ref{thm:scalar-constant-characteristics}) that it cannot be smooth. For this type of initial data we will have to develop a theory of discontinuous solutions, or "shock waves."

\begin{figure}[H]
\centering
\includegraphics[width=.78\linewidth]{figures/characteristic-gap-source.png}
\caption{Characteristics do not specify the solution in the blank region.}\label{fig:characteristic-gap}
\end{figure}

\begin{figure}[H]
\centering
\includegraphics[width=.78\linewidth]{figures/characteristic-overlap-source.png}
\caption{Characteristics overlap.}\label{fig:characteristic-overlap}
\end{figure}

Note that smoothing out the data does not help matters in this case; it merely delays the problem. In fact, the following theorem shows that the problem of overlapping characteristics and the development of singularities is a generic problem.

\begin{figure}[H]
\centering
\includegraphics[width=.78\linewidth]{figures/continuous-characteristic-crossing-source.png}
\caption{Intersecting characteristics from continuous initial data.}\label{fig:continuous-characteristic-crossing}
\end{figure}

\begin{Theorem}
\label{thm:scalar-finite-time-blowup}
If $f''>0$ and the initial data $u_0$ is not monotone increasing, then the Cauchy problem for \eqref{eq:scalar-conservation-law-quasilinear} does not have a $C^1$ solution defined on the entire upper half-plane $(x,t)\in(-\infty,\infty)\times[0,\infty)$.
\end{Theorem}

\begin{Proof}
The proof simply depends on the observation that if $f'(u_0(x_1))>f'(u_0(x_2))$, for $x_1<x_2$, the characteristics emanating from $x_1$ and $x_2$ will intersect in finite time (cf. Figure~\ref{fig:continuous-characteristic-crossing}).
\end{Proof}

\subsubsection{The $p$ System}\label{subsec:p-system-blowup}

We use more analytical techniques to prove blowup for the $p$ system, where characteristic curves are no longer so simple. We make use of the diagonal form of the system \eqref{eq:p-system-diagonal-first}, \eqref{eq:p-system-diagonal-second} given by changing to Riemann invariant coordinates in state space. Since Theorem~\ref{thm:p-system-characteristic-invariants} implies that $\rho_1$ and $\rho_2$ are constant along their respective characteristic curves, we cannot expect them to become unbounded as long as the solution stays $C^1$. However, if we examine the evolution of the slopes $\rho_{1,x}$ and $\rho_{2,x}$, we can expect something to go wrong. Thus, we differentiate \eqref{eq:p-system-diagonal-first} and \eqref{eq:p-system-diagonal-second} with respect to $x$ to obtain
\begin{equation}\label{eq:p-system-first-slope-equation}
\rho_{1,xt}+\hat\lambda(\rho_1-\rho_2)\rho_{1,xx}=-\hat\lambda'(\rho_1-\rho_2)(\rho_{1,x}^2-\rho_{1,x}\rho_{2,x}),
\end{equation}

\begin{equation}\label{eq:p-system-second-slope-equation}
\rho_{2,xt}-\hat\lambda(\rho_1-\rho_2)\rho_{2,xx}=-\hat\lambda'(\rho_1-\rho_2)(\rho_{2,x}^2-\rho_{1,x}\rho_{2,x}).
\end{equation}
The product terms $\rho_{1,x}\rho_{2,x}$ cause an inconvenient coupling, but we can get rid of them by using the change of variables $r:=\hat\lambda^{1/2}\rho_{1,x}$, $s:=\hat\lambda^{1/2}\rho_{2,x}$.

Under this change our system becomes
\begin{equation}\label{eq:p-system-weighted-first-slope}
r_t+\hat\lambda(\rho_1-\rho_2)r_x=-\hat\lambda^{-1/2}(\rho_1-\rho_2)\hat\lambda'(\rho_1-\rho_2)r^2,
\end{equation}

\begin{equation}\label{eq:p-system-weighted-second-slope}
s_t-\hat\lambda(\rho_1-\rho_2)s_x=-\hat\lambda^{-1/2}(\rho_1-\rho_2)\hat\lambda'(\rho_1-\rho_2)s^2.
\end{equation}
Hence, the derivatives of $r$ and $s$ along characteristics are proportional to $r^2$ and $s^2$, respectively. With this in mind, consider the following lemma.

\begin{Lemma}
\label{lem:riccati-finite-time-blowup}
Let $z$ be the solution of the ODE initial-value problem
\begin{equation}\label{eq:riccati-positive-coefficient}
z'(t)=a(t)z^2(t),
\end{equation}

\begin{equation}\label{eq:riccati-positive-initial-value}
z(0)=m>0,
\end{equation}
where $0<B\leq a(t)\leq A$. Then $z$ becomes infinite (has a vertical asymptote) at time $t_c\in((mA)^{-1},(mB)^{-1})$. Furthermore, for $t<t_c$ we have the estimate
\begin{equation}\label{eq:riccati-blowup-bounds}
\dfrac1{A(t_c-t)}<z(t)<\dfrac1{B(t_c-t)}.
\end{equation}

\end{Lemma}

Using this and the calculations above, we can show

\begin{Theorem}
\label{thm:p-system-finite-time-blowup}
Assume that
\[
\sup_{x_1,x_2\in\mathbb{R}}\hat\lambda'(\rho_1(x_1,0)-\rho_2(x_2,0))<0
\]
and that
\[
m:=\max[\sup_xr(x,0),\sup_xs(x,0)]>0.
\]
Let
\[
\begin{aligned}
A&:=-\inf_{x_1,x_2\in\mathbb{R}}\hat\lambda^{-1/2}\hat\lambda'(\rho_1(x_1,0)-\rho_2(x_2,0)),\\
B&:=-\sup_{x_1,x_2\in\mathbb{R}}\hat\lambda^{-1/2}\hat\lambda'(\rho_1(x_1,0)-\rho_2(x_2,0)).
\end{aligned}
\]
Then at least one of $r$ and $s$ becomes unbounded at a time between $(mA)^{-1}$ and $(mB)^{-1}$.
\end{Theorem}

The proof of this theorem and Lemma~\ref{lem:riccati-finite-time-blowup} are left to the reader as exercises.
